\chapter{Introdução}

A onipresença da criptografia na infraestrutura digital moderna (internet, sistemas bancários, IoT e comunicações seguras).

\section{Contextualização do Tema}

Textos 

\section{Justificativa/Problema de Pesquisa}

O conflito entre o determinismo computacional e a necessidade de aleatoriedade em~sistemas criptográficos.

\section{Motivação/Hipotese e Ameaca Tecnológica}

O impacto da computação quântica --- especialmente por meio do Algoritmo de Shor --- sobre os modelos atuais de segurança e pseudoaleatoriedade.

\section{Objetivos}

Objetivo Geral: Analisar a evolução dos geradores pseudoaleatórios criptograficamente seguros frente à transição para a criptografia pós-quântica.
Objetivos Específicos:
Estudar PRNGs e CSPRNGs clássicos
Analisar o impacto dos algoritmos quânticos
Investigar os novos paradigmas de pseudoaleatoriedade baseados em reticulados
Avaliar a adoção desses conceitos no ecossistema Java


\section{Metodologia}

Pesquisa bibliográfica, análise documental de padrões criptográficos e estudo prático no ecossistema Java.